\exnumber{02}
\extitle{UI I - Ex 3 - Color combinations}
\exfiles{ex02/Figma file, named Name\_FirstName\_DeliverableName\_Date\_VersionNumber}
\exauthorize{Those named in the exercise}
\makeheaderfiles

\textbf{Instructions:}

Let's browse!

Find 3 digital products (websites or apps) that use each kind of color palette combination:

\begin{itemize}
\item Monochromatic
\item Analogous
\item Complementary
\item Split complementary
\item Triadic
\item Double complementary
\end{itemize}

\textbf{Deliverable:}

\begin{itemize}
\item A Figma file including for each digital product (websites or apps)
\begin{itemize}
\item screenshots of the main page
\item the matching URL
\item the main colors and
\item the kind of color combination.
\end{itemize}
\item Upload the Figma file on Github, labeled as follows: "Name\_FirstName\_DeliverableName\_Date\_VersionNumber". For example, it can be labeled as "Doe\_John\_MidFiPrototype\_05242024\_V1".
\item Add a title within the file.
\item Pay attention to the global aspect of the file. It must be clear, simple and easy to read. You can get inspiration on canva.com but don't overload your design with too much details!
\end{itemize}

\textbf{Tips:}

\begin{itemize}
\item You can find HSB (Hue, Saturation, Brightness) as well as HSV (Hue, Saturation, Value) depending of the source.
\end{itemize}

\textbf{Resources:}

\begin{itemize}
\item https://coloursandmaterials.wordpress.com/2014/09/16/colour-system/ \url{https://coloursandmaterials.wordpress.com/2014/09/16/colour-system/}
\item Color wheel combinations
\end{itemize}

\textbf{Monochromatic examples:}

\begin{itemize}
\item https://www.benefitcosmetics.com/ \url{https://www.benefitcosmetics.com/} 
\item http://www.jwilouvres.com.au/ \url{http://www.jwilouvres.com.au/} 
\item http://blank.com.pt/ \url{http://blank.com.pt/}
\end{itemize}

\textbf{Analogous examples:}

\begin{itemize}
\item https://www.natuurenmilieu.nl/ \url{https://www.natuurenmilieu.nl/} 
\item https://campesinorum.com/ \url{https://campesinorum.com/} 
\item https://www.lobster.es/ \url{https://www.lobster.es/}
\end{itemize}

\textbf{Complementary examples:}

\begin{itemize}
\item https://www.rolandgarros.com/ \url{https://www.rolandgarros.com/fr-fr} 
\item https://squilla.io/ \url{https://squilla.io/} 
\item https://gomim.com/en/ \url{https://gomim.com/en/}
\end{itemize}

\textbf{Split complementary examples:}

\begin{itemize}
\item https://flythenest.io/en/ \url{https://flythenest.io/en/} 
\item https://www.kineuphorics.com/ \url{https://www.kineuphorics.com/} 
\item https://www.ma-tea.de/ \url{https://www.ma-tea.de/}
\end{itemize}

\textbf{Triadic examples}

\begin{itemize}
\item https://normalnow.com/ \url{https://normalnow.com/} 
\item https://brainbakery.com/en/ \url{https://brainbakery.com/en/} 
\item https://www.barkbox.com/ \url{https://www.barkbox.com/}
\end{itemize}

\textbf{Dual complementary or Tetradic examples}

\begin{itemize}
\item https://www.airtable.com/ \url{https://www.airtable.com/} 
\item https://calendar.google.com/ \url{https://calendar.google.com/} 
\item https://www.ebay.com/ \url{https://www.ebay.com/}
\end{itemize}

